\chapter*{Abstract}
\addcontentsline{toc}{chapter}{Abstract}

The identification of patient symptoms in time and correctly to enhance access to healthcare and decrease avoidable medical visits is relevant. I have compared five machine learning classifiers in this thesis (Naive Bayes, Logistic Regression, Support Vector Machine, Random Forest, and XGBoost) to predict diseases with a number of classes using binary symptom vectors. I compare them on two publicly available Kaggle datasets, a small dataset consisting of about 5,000 records with 132 symptoms and 41 diseases and a larger dataset consisting of about 246,000 records with a considerably larger set of classes. In the case of Dataset 1, I use stratified cross-validation using grid search. In Dataset 2, I apply fixed hyperparameters to make the experiments computationally feasible. I compare all models based on accuracy, precision, recall and F1-score to determine the approaches that are best to use in predicting symptoms. Limitations of the study are the lack of external clinical validation, binary representation of symptoms, and simplified one-disease formulation. Besides the metric findings, I address real-world trade-offs of interpretability, computational cost, and regulatory requirements on the EU AI Act and GDPR to actual healthcare booking systems.

\vspace{0.5cm}
\noindent\textbf{Keywords:} machine learning, disease prediction, symptom-based classification, healthcare, Random Forest, XGBoost, comparative study
